\chapter{Exchanges, Brokers and Venues}\label{ch:m1:exchanges-brokers-venues}

The company that owns the New York Stock \omterm{def:m1:exchanges-brokers-venues:exchange}{Exchange} earned more in 2025 from
selling data and network connections at its \omterm{def:m1:exchanges-brokers-venues:exchange}{exchanges} than from every share
and every stock option traded on them. The famous floor is a small part of a
business that lists companies, matches orders, clears trades, rents cabinets
in a data centre and licenses the record of what happened. A trading firm is
a customer of all five --- and cannot send a single order before deciding
through whom, and under whose name, it reaches the matching engine. This
chapter describes what an \omterm{def:m1:exchanges-brokers-venues:exchange}{exchange} sells, the chain of intermediaries
between a trader and the match, and the zoo of venues that are not
\omterm{def:m1:exchanges-brokers-venues:exchange}{exchanges}.

\section{What an exchange sells}

\begin{definition}[Order]\label{def:m1:exchanges-brokers-venues:order}
An \emph{order}\index{order} is a binding instruction to buy or to sell a
stated quantity of an instrument, carrying at least a side, a quantity, a
price condition (a limit, or ``at market'') and a validity. It remains the
responsibility of whoever sent it until it is executed, cancelled or expired.
\end{definition}

\begin{definition}[Exchange]\label{def:m1:exchanges-brokers-venues:exchange}
An \emph{exchange}\index{exchange} is a regulated organisation that admits
instruments to trading, admits members, and operates a system in which
members' orders interact under public, non-discretionary rules to form
trades and prices. It supervises its own market and is itself supervised by
a public authority.
\end{definition}

The definition hides five products, sold to different customers.

\begin{method}[Reading an exchange group's revenue]\label{met:m1:exchanges-brokers-venues:products}
Sort each revenue line into:
\begin{enumerate}
\item \textbf{Listing} --- paid by issuers, annually, for admission to
  trading.
\item \textbf{Transaction} --- paid by members per share, contract or unit of
  value matched, net of any rebates paid back to them.
\item \textbf{Clearing} --- where the group owns the clearing house
  (\cref{ch:m1:clearing-and-settlement}), paid per trade cleared and earned
  on collateral held.
\item \textbf{Data} --- licences for real-time and historical prices, paid by
  everyone who looks at or computes on them.
\item \textbf{Access} --- ports, cross-connects, cabinets and power in the
  \omterm{def:m1:exchanges-brokers-venues:exchange}{exchange}'s data centre, paid by whoever needs to be close
  (One Quant Book 14).
\end{enumerate}
Lines 1, 4 and 5 recur whatever volumes do; 2 and 3 move with the market.
Divide 2 and 3 by the volume to obtain the \emph{revenue per contract} or
per share: the price of the \omterm{def:m1:exchanges-brokers-venues:exchange}{exchange}'s core service.
\end{method}

\begin{dated}{2026-09}{Two exchange groups in their own numbers}\label{dat:m1:exchanges-brokers-venues:groups}
\textbf{Intercontinental \omterm{def:m1:exchanges-brokers-venues:exchange}{Exchange}}, owner of the New York Stock \omterm{def:m1:exchanges-brokers-venues:exchange}{Exchange},
reported 2025 net revenues of \$9.9 billion, of which \$5\,411 million in its
\omterm{def:m1:exchanges-brokers-venues:exchange}{Exchanges} segment: energy \$2\,182 million, agricultural products and metals
\$233 million, financial futures \$608 million, cash equities and equity
options \$467 million, over-the-counter and other \$395 million, data and
connectivity services \$1\,031 million, listings \$495 million.
\textbf{CME Group} reported 2025 revenue of \$6.5 billion, of which
\$5\,281 million in clearing and transaction fees and \$803 million in
market data; its average rate per contract was \$0.696.
\end{dated}

\begin{omfigure}
\begin{tikzpicture}
\begin{axis}[width=11.5cm,height=5.2cm,xbar,bar width=7pt,
  xlabel={2025 net revenues (\$ million)},
  ytick=data,yticklabels from table={figdata/markets-1/04-exchanges-brokers-venues/ice_exchanges_2025.csv}{label},
  yticklabel style={font=\small},tick label style={font=\small},label style={font=\small},
  y dir=reverse,xmin=0,xmax=2500,xtick={0,500,1000,1500,2000,2500},scaled x ticks=false,
  xticklabel style={/pgf/number format/1000 sep={\,}},enlarge y limits=0.1,grid=major,grid style={gray!25},
  table/col sep=comma]
\addplot[fill=omDef!60,draw=omDef] table[x=usd_m,y=k]{figdata/markets-1/04-exchanges-brokers-venues/ice_exchanges_2025.csv};
\end{axis}
\end{tikzpicture}

\omcaption{The \omterm{def:m1:exchanges-brokers-venues:exchange}{Exchanges} segment of one group in 2025. Data and listings ---
the two bottom bars, 28\% of the segment --- do not depend on volumes; equity
trading is under a tenth. Data: the group's earnings release, as in
\cref{dat:m1:exchanges-brokers-venues:groups}.}\label{fig:m1:exchanges-brokers-venues:ice}
\end{omfigure}

\begin{example}[What a futures contract pays its exchange]\label{ex:m1:exchanges-brokers-venues:rpc}
At \$0.696 a contract, \$5\,281 million of clearing and transaction fees
correspond to about 7.6 billion contract sides charged in the year, some 30
million per trading day. A single equity index contract controls several
hundred thousand dollars of stock: the \omterm{def:m1:exchanges-brokers-venues:exchange}{exchange}'s fee is a small fraction of
a basis point of the value traded. \omterm{def:m1:exchanges-brokers-venues:exchange}{Exchanges} are, like \omterm{def:m1:what-a-trading-firm-does:market-maker}{market makers}, a
business of tiny numbers multiplied by enormous ones --- with the difference
that an \omterm{def:m1:exchanges-brokers-venues:exchange}{exchange}'s customers cannot easily take a contract elsewhere.
\end{example}

\begin{remark}[Why futures exchanges earn more than stock exchanges]\label{rem:m1:exchanges-brokers-venues:silo}
A share is the same share on every venue: US and European law force
competition between stock \omterm{def:m1:exchanges-brokers-venues:exchange}{exchanges}, and matching fees were competed down to
almost nothing. A futures contract is a creature of the \omterm{def:m1:exchanges-brokers-venues:exchange}{exchange} that lists
it and clears at that \omterm{def:m1:exchanges-brokers-venues:exchange}{exchange}'s clearing house; a position opened there can
only be closed there. The ``vertical silo'' of listing, matching and clearing
is why \cref{fig:m1:exchanges-brokers-venues:ice} looks the way it does.
\end{remark}

\section{Members, brokers and access}

\begin{definition}[Exchange member]\label{def:m1:exchanges-brokers-venues:member}
An \emph{exchange member}\index{exchange member} (or \emph{participant}) is a
firm admitted by the \omterm{def:m1:exchanges-brokers-venues:exchange}{exchange} to send orders to its system in its own name.
It must be authorised by a regulator, meet capital requirements, accept the
rulebook, and have an arrangement to clear its trades.
\end{definition}

\begin{definition}[Broker]\label{def:m1:exchanges-brokers-venues:broker}
A \emph{broker}\index{broker} is a member that sends orders on behalf of
clients who are not members. The \omterm{def:m1:exchanges-brokers-venues:exchange}{exchange} knows only the broker: towards the
market the broker is responsible for every order, and for paying for every
trade, of every client behind it.
\end{definition}

That last sentence explains the structure of the industry. Because the
\omterm{def:m1:exchanges-brokers-venues:broker}{broker} answers for its clients, it must control what they send; because the
controls take time, the fastest firms want to avoid them; and regulators
decide how far they may.

\begin{definition}[Direct market access and sponsored access]\label{def:m1:exchanges-brokers-venues:dma}
Under \emph{direct market access}\index{direct market access} (DMA) a client
sends orders through the \omterm{def:m1:exchanges-brokers-venues:broker}{broker}'s systems, which check them and pass them to
the \omterm{def:m1:exchanges-brokers-venues:exchange}{exchange} under the \omterm{def:m1:exchanges-brokers-venues:broker}{broker}'s membership. Under
\emph{sponsored access}\index{sponsored access} the client connects to the
\omterm{def:m1:exchanges-brokers-venues:exchange}{exchange} directly, with its own lines and machines, using the \omterm{def:m1:exchanges-brokers-venues:broker}{broker}'s
membership identifier; the \omterm{def:m1:exchanges-brokers-venues:broker}{broker}'s risk checks are applied to that flow
without its passing through the \omterm{def:m1:exchanges-brokers-venues:broker}{broker}'s order-handling infrastructure.
\end{definition}

\begin{omfigure}
\begin{tikzpicture}[font=\small,
  box/.style={draw=omDef,rounded corners=2pt,minimum width=2.4cm,minimum height=0.75cm,align=center,fill=omDef!4},
  chk/.style={draw=omThm,rounded corners=2pt,minimum width=2.2cm,minimum height=0.75cm,align=center,fill=omThm!5},
  ex/.style={draw=omExo,rounded corners=2pt,minimum width=1.6cm,minimum height=5.9cm,align=center,fill=omExo!6},
  flow/.style={-{Stealth[length=2mm]},thick,omDef}]
\node[ex] (ex) at (13.2,-2.55) {Exchange};
\foreach \y/\lab in {0/Broker's algorithm,-1.7/Direct market access,-3.4/Sponsored access,-5.1/Own membership}{
  \node[anchor=west,font=\small\bfseries] at (-0.9,\y+0.74) {\lab};
}
\node[box] (c1) at (0.4,0) {Client};
\node[box] (b1) at (4.2,0) {Broker's trader\\and algorithm};
\node[chk] (k1) at (8.4,0) {Broker's\\risk checks};
\draw[flow] (c1) -- (b1); \draw[flow] (b1) -- (k1); \draw[flow] (k1) -- (k1 -| ex.west);
\node[box] (c2) at (0.4,-1.7) {Client's\\algorithm};
\node[chk] (k2) at (8.4,-1.7) {Broker's\\risk checks};
\draw[flow] (c2) -- node[above,font=\footnotesize]{through the broker's systems} (k2); \draw[flow] (k2) -- (k2 -| ex.west);
\node[box] (c3) at (0.4,-3.4) {Client's\\algorithm};
\node[chk] (k3) at (8.4,-3.4) {Broker-controlled\\checks at the client};
\draw[flow] (c3) -- node[above,font=\footnotesize]{client's own lines} (k3); \draw[flow] (k3) -- (k3 -| ex.west);
\node[box] (c4) at (0.4,-5.1) {Member's\\algorithm};
\node[chk] (k4) at (8.4,-5.1) {Member's own\\risk checks};
\draw[flow] (c4) -- (k4); \draw[flow] (k4) -- (k4 -| ex.west);
\end{tikzpicture}

\omcaption{Four ways to reach the same matching engine, from slowest and
cheapest to set up (top) to fastest and most demanding (bottom). In every
row a regulated firm's pre-trade checks (red) sit between the algorithm and
the \omterm{def:m1:exchanges-brokers-venues:exchange}{exchange}: what changes is whose they are and where they
run.}\label{fig:m1:exchanges-brokers-venues:access}
\end{omfigure}

\begin{remark}[The end of naked access]\label{rem:m1:exchanges-brokers-venues:15c35}
Until 2010 some US \omterm{def:m1:exchanges-brokers-venues:broker}{brokers} lent their identifier to fast clients with no
pre-trade control at all (``naked'' access). On 3 November 2010 the SEC
adopted Rule 15c3-5, the \emph{market access rule}: a \omterm{def:m1:the-sell-side:sell-side}{broker-dealer} with
market access, or providing it, must apply pre-trade risk controls ---
credit and capital thresholds, rejection of erroneous orders, regulatory
checks --- that are under its \emph{direct and exclusive control}. European
rules on algorithmic trading and direct electronic access impose the
equivalent. The checks a firm builds to satisfy them are a component of the
miniature firm of this series (One Quant Book 13); what they must catch is
the lesson of the 2012 incident studied there.
\end{remark}

\begin{method}[Choosing an access model]\label{met:m1:exchanges-brokers-venues:choose}
\begin{enumerate}
\item \textbf{Latency need.} If the strategy competes on speed at the
  matching engine, only \omterm{def:m1:exchanges-brokers-venues:dma}{sponsored access} or membership will do.
\item \textbf{Volume.} Compute the all-in monthly cost
  $F + \varphi V$ of each model (fixed cost $F$, per-share charge $\varphi$,
  volume $V$) and find the crossover volumes $V_{ab} = (F_b - F_a)/(\varphi_a
  - \varphi_b)$.
\item \textbf{Regulatory appetite.} Membership means becoming a regulated
  \omterm{def:m1:the-sell-side:sell-side}{broker-dealer} or investment firm: capital, compliance staff, reporting,
  examinations.
\item \textbf{Counterparty.} A \omterm{def:m1:exchanges-brokers-venues:broker}{broker} can cut a client off in a crisis; a
  member depends only on its clearing firm (\cref{ch:m1:clearing-and-settlement}).
\end{enumerate}
The decision is revisited as the firm grows;
\cref{ch:m1:getting-access-equities,ch:m1:getting-access-futures-options}
work through real fee schedules.
\end{method}

\begin{example}[Four models, three crossovers]\label{ex:m1:exchanges-brokers-venues:crossover}
Take illustrative charges of 0.30 cent a share with no fixed cost for a
\omterm{def:m1:exchanges-brokers-venues:broker}{broker}'s algorithm; 0.10 cent plus \$5\,000 a month for DMA; 0.04 cent plus
\$25\,000 a month for \omterm{def:m1:exchanges-brokers-venues:dma}{sponsored access} (the client now pays for its own
lines and colocation); and \$150\,000 a month with no per-share charge for a
firm's own membership (compliance, capital, memberships). The crossovers are
$5\,000/0.002 = 2.5$ million shares a month, $20\,000/0.0006 = 33$ million,
and $125\,000/0.0004 = 313$ million. A firm trading 15 million shares a
\emph{day} is past the last one; a fund trading one million a month should
not even consider DMA (\cref{fig:m1:exchanges-brokers-venues:cost}).
\end{example}

\begin{omfigure}
\begin{tikzpicture}
\begin{axis}[width=11.5cm,height=5.4cm,xmode=log,ymode=log,
  xlabel={shares traded per month (million)},ylabel={all-in cost (cents per share)},
  tick label style={font=\small},label style={font=\small},
  xmin=0.1,xmax=2000,ymin=0.005,ymax=20,grid=major,grid style={gray!25},
  legend columns=4,legend style={font=\footnotesize,at={(0.5,-0.32)},anchor=north,draw=none,fill=none,/tikz/every even column/.append style={column sep=6pt}},
  table/col sep=comma]
\addplot[omExo,very thick] table[x=shares_m,y=broker_algorithm]{figdata/markets-1/04-exchanges-brokers-venues/access_cost.csv};
\addplot[omMeth,very thick] table[x=shares_m,y=direct_market_access]{figdata/markets-1/04-exchanges-brokers-venues/access_cost.csv};
\addplot[omDef,very thick] table[x=shares_m,y=sponsored_access]{figdata/markets-1/04-exchanges-brokers-venues/access_cost.csv};
\addplot[omThm,very thick] table[x=shares_m,y=own_membership]{figdata/markets-1/04-exchanges-brokers-venues/access_cost.csv};
\legend{broker algorithm,DMA,sponsored access,own membership}
\end{axis}
\end{tikzpicture}

\omcaption{Cost per share of four access models against monthly volume (both
axes logarithmic). The cheapest model changes three times over four decades
of volume. Data: the illustrative charges of
\cref{ex:m1:exchanges-brokers-venues:crossover}, computed by the chapter's
script.}\label{fig:m1:exchanges-brokers-venues:cost}
\end{omfigure}

\section{Venues that are not exchanges}

Matching orders under rules does not require an \omterm{def:m1:exchanges-brokers-venues:exchange}{exchange} licence.

\begin{definition}[Alternative trading system]\label{def:m1:exchanges-brokers-venues:ats}
In US law an \emph{alternative trading system}\index{alternative trading system}
(ATS) is a system that brings together buyers and sellers of securities as an
\omterm{def:m1:exchanges-brokers-venues:exchange}{exchange} does but is operated by a \omterm{def:m1:the-sell-side:sell-side}{broker-dealer} under Regulation ATS
(adopted in December 1998) instead of registering as an \omterm{def:m1:exchanges-brokers-venues:exchange}{exchange}. It does not
list securities and does not regulate its participants.
\end{definition}

\begin{definition}[Multilateral trading facility]\label{def:m1:exchanges-brokers-venues:mtf}
In EU law a \emph{multilateral trading facility}\index{multilateral trading facility}
(MTF) is a multilateral system, operated by an investment firm or a market
operator, which brings together multiple third-party buying and selling
interests in financial instruments, under non-discretionary rules, in a way
that results in a contract.
\end{definition}

Most \emph{dark pools} (\cref{ch:m1:us-equity-market-structure}) are ATSs;
the large pan-European stock venues that compete with the national \omterm{def:m1:exchanges-brokers-venues:exchange}{exchanges}
are MTFs (\cref{ch:m1:european-equity-market-structure}). Both trade
instruments that an \omterm{def:m1:exchanges-brokers-venues:exchange}{exchange} listed, free-riding on its listing and
supervision work, which is one reason \omterm{def:m1:exchanges-brokers-venues:exchange}{exchanges} charge what they do for
data.

\begin{proposition}[Break-even share of a new venue]\label{prop:m1:exchanges-brokers-venues:breakeven}
A venue with annual fixed costs $C$ enters a market trading $M$ shares a day
over $D$ days a year. It keeps a net fee $\nu$ per share matched and earns
market-data revenue in proportion to its share, $\delta$ at 100\%. Its
profit at market share $x$ is $x\,(M D \nu + \delta) - C$ and its break-even
share is
\[
x^\star \;=\; \frac{C}{M D \nu + \delta}.
\]
\end{proposition}
\begin{proof}
Matched volume is $xMD$ shares a year, each earning $\nu$; data revenue is
$x\delta$; costs are fixed.
\end{proof}

\begin{example}[Five percent or nothing]\label{ex:m1:exchanges-brokers-venues:five}
With $M = 10$ billion shares a day, $D = 252$, $\nu = 0.03$ cent, $\delta =
\$50$ million and $C = \$40$ million, $x^\star = 40/(756+50) = 5.0\%$. At a
10\% share the venue earns \$40.6 million; at 2\% it loses \$23.9 million.
But share is not a free parameter: traders send orders where they expect to
be filled, that is, where the other orders already are. A new venue must
\emph{buy} its first percent --- with rebates above its fees, with equity
stakes offered to the \omterm{def:m1:exchanges-brokers-venues:broker}{brokers} who route to it, or with a rule that some group
of traders values (a speed bump, a midpoint book). Liquidity attracts
liquidity; the history of venues is a short list of survivors.
\end{example}

\begin{omfigure}
\begin{tikzpicture}
\begin{axis}[width=10.5cm,height=4.8cm,
  xlabel={market share of the new venue (\%)},ylabel={annual profit (\$ million)},
  tick label style={font=\small},label style={font=\small},
  xmin=0,xmax=20,ymin=-45,ymax=125,grid=major,grid style={gray!25},table/col sep=comma]
\addplot[omDef,very thick] table[x=share_pct,y=profit_m]{figdata/markets-1/04-exchanges-brokers-venues/venue_profit.csv};
\addplot[omThm,dashed] coordinates {(4.96,-45) (4.96,125)};
\node[omThm,font=\small,anchor=west] at (axis cs:5.3,100) {$x^\star = 5.0\%$};
\end{axis}
\end{tikzpicture}

\omcaption{Profit of a new venue against its market share, with the
parameters of \cref{ex:m1:exchanges-brokers-venues:five}. Below five percent
it burns its whole cost base. Data: computed by the chapter's script.}
\end{omfigure}

\section{The consolidated view}

\begin{definition}[Market data feed]\label{def:m1:exchanges-brokers-venues:feed}
A \emph{market data feed}\index{market data feed} is the stream of messages
by which a venue publishes its orders, quotes and trades. A
\emph{consolidated} feed merges the feeds of all venues trading the same
instruments into one best bid, one best offer and one tape of trades.
\end{definition}

Once one share trades in twenty places, three problems appear, and each
gets its own chapter. \emph{What is the price?} Somebody must merge the
feeds, and the merged view is always a little older than its sources
(\cref{ch:m1:us-equity-market-structure,ch:m1:reading-market-data}).
\emph{Where should an order go?} A \omterm{def:m1:exchanges-brokers-venues:broker}{broker}'s duty of best execution becomes a
routing problem solved by software (One Quant Book 10). \emph{Who supervises
a manipulation that uses three venues at once?} Cross-market surveillance is
the regulators' answer, and its detectors are in One Quant Book 9.

\begin{remark}[Every venue has a code]\label{rem:m1:exchanges-brokers-venues:mic}
Systems identify venues by the four-character \emph{market identifier code}
of ISO~10383: \texttt{XNYS} is the New York Stock \omterm{def:m1:exchanges-brokers-venues:exchange}{Exchange}, \texttt{XNAS}
Nasdaq. An \emph{operating} code names the organisation and \emph{segment}
codes name its individual markets; the registry is published monthly and is
free. The code, not the name, is what appears in trade reports, routing
tables and this chapter's build.
\end{remark}

\section{Tutorial: reading an exchange's accounts}\label{tut:m1:exchanges-brokers-venues:accounts}

\begin{tutorial}
\textbf{Goal.} Turn two earnings releases into the quantities a trading firm
cares about, then compute the access crossovers. \textbf{End state:} the
shares quoted in \cref{fig:m1:exchanges-brokers-venues:ice} and the three
crossovers of \cref{ex:m1:exchanges-brokers-venues:crossover}.
\begin{tutsteps}
\item \textbf{Enter the published lines}, with their source in a comment ---
a number without a source is a rumour.
\omcode{code/markets-1/04-exchanges-brokers-venues/python/venue_econ.py}{4}{15}{Two groups' 2025 revenue lines as data.}\label{lst:m1:exchanges-brokers-venues:data}
\item \textbf{Compute shares.} Recurring lines (data, connectivity,
listings) are $1\,526/5\,411 = 28\%$ of the segment; cash equities and
equity options $8.6\%$; energy $40\%$. At CME, clearing and transaction fees
are $81\%$ of revenue and market data $12\%$.
\item \textbf{Model the access decision.}
\omcode{code/markets-1/04-exchanges-brokers-venues/python/venue_econ.py}{23}{48}{Access models as data, their monthly cost and the crossover volume.}\label{lst:m1:exchanges-brokers-venues:access}
\item \textbf{Find the cheapest model} at 1, 10, 100 and 1\,000 million
shares a month: \omterm{def:m1:exchanges-brokers-venues:broker}{broker} algorithm, DMA, \omterm{def:m1:exchanges-brokers-venues:dma}{sponsored access}, own membership.
\end{tutsteps}
\textbf{What to change next.} Add a latency column and exclude models slower
than a threshold: the cheapest \emph{admissible} model at 5 million shares a
month changes. Then add the \omterm{def:m1:exchanges-brokers-venues:exchange}{exchange}'s own tiered fees
(\cref{ch:m1:getting-access-equities}), which depend on volume too.
\end{tutorial}

\section{Build: the venue registry}\label{bld:m1:exchanges-brokers-venues:registry}

\begin{build}
\bfield{Purpose} Every order, fill, fee and market-data message in the
miniature firm names a venue. The registry is the one place that knows what
each venue is.
\bfield{Interface} \texttt{Venue(mic, operating\_mic, name, kind, country,
currency, fee\_model)} with \texttt{kind} in \texttt{exchange}, \texttt{ats},
\texttt{mtf}, \texttt{si}, \texttt{otc}; and
\texttt{Registry.load(csv\_path)}, \texttt{get(mic)},
\texttt{by\_kind(kind)}, \texttt{segments\_of(operating\_mic)}.
\bfield{Rules} Codes are four upper-case alphanumerics; a duplicate code, an
unknown kind or a segment whose operating code is absent is rejected at load
time, not at first use. The registry is immutable after loading.
\bfield{Acceptance tests} \texttt{code/firm/venues/tests/}, run against the
sample file \texttt{data/markets-1/venues\_sample.csv}.
\bfield{Stretch} Load the full ISO~10383 file from the registration
authority and report how many active segment codes each operating code has.
\end{build}

\begin{omsources}
\item Intercontinental Exchange, \emph{Intercontinental Exchange Reports
Strong Full Year 2025 Results}, 5 February 2026.
\item CME Group, \emph{CME Group Inc. Reports Fourth Consecutive Year of
Record Annual Revenue}, 4 February 2026.
\item US Securities and Exchange Commission, \emph{SEC Adopts New Rule
Preventing Unfiltered Market Access}, press release 2010-210, and Release
34-63241 (Rule 15c3-5).
\item 17 CFR 242.300 (Regulation ATS, definitions); Directive 2014/65/EU
(MiFID II), article 4(1)(22).
\item ISO 10383 Registration Authority, \emph{Market identifier codes}.
\item L.~Harris, \emph{Trading and Exchanges}, Oxford University Press, 2003,
chapters 3, 7 and 26.
\end{omsources}

\section{Exercises}

\begin{exercise}[$\star$]\label{exo:m1:exchanges-brokers-venues:1}
From \cref{dat:m1:exchanges-brokers-venues:groups}, compute for the ICE
\omterm{def:m1:exchanges-brokers-venues:exchange}{Exchanges} segment the share of each of: energy; cash equities and equity
options; data and connectivity; listings.
\end{exercise}

\begin{exercise}[$\star$]\label{exo:m1:exchanges-brokers-venues:2}
A futures contract has a notional value of \$280\,000 and the \omterm{def:m1:exchanges-brokers-venues:exchange}{exchange} earns
\$0.70 per side. Express the fee of a round trip (buy then sell) in basis
points of the notional.
\end{exercise}

\begin{exercise}[$\star$]\label{exo:m1:exchanges-brokers-venues:3}
For each situation name the access model and say whose pre-trade risk checks
apply: (a) a fund manager phones an order to a \omterm{def:m1:the-sell-side:sales-trader}{sales-trader}; (b) a
quantitative fund's algorithm sends orders through its \omterm{def:m1:exchanges-brokers-venues:broker}{broker}'s gateway;
(c) a trading firm colocated at the \omterm{def:m1:exchanges-brokers-venues:exchange}{exchange} uses its \omterm{def:m1:exchanges-brokers-venues:broker}{broker}'s identifier;
(d) a \omterm{def:m1:what-a-trading-firm-does:market-maker}{market maker} registered as a \omterm{def:m1:the-sell-side:sell-side}{broker-dealer}.
\end{exercise}

\begin{exercise}[$\star\star$]\label{exo:m1:exchanges-brokers-venues:4}
A firm trades 60 million shares a month. With the charges of
\cref{ex:m1:exchanges-brokers-venues:crossover}, compute the monthly cost of
each model and the saving of the best over the second best. At what monthly
volume does it become worth applying for membership?
\end{exercise}

\begin{exercise}[$\star\star$]\label{exo:m1:exchanges-brokers-venues:5}
A sponsored-access provider raises its per-share charge from 0.04 to 0.06
cent. Recompute the two crossovers that involve \omterm{def:m1:exchanges-brokers-venues:dma}{sponsored access}. Which kind
of client does the increase push away?
\end{exercise}

\begin{exercise}[$\star\star$]\label{exo:m1:exchanges-brokers-venues:6}
With the parameters of \cref{ex:m1:exchanges-brokers-venues:five}, the
incumbents cut their fees and the new venue must follow: $\nu$ falls to 0.02
cent. Compute the new break-even share and the profit at 10\%.
\end{exercise}

\begin{exercise}[$\star\star\star$]\label{exo:m1:exchanges-brokers-venues:7}
\emph{Coding.} Write \texttt{cheapest\_admissible(volume, max\_latency)}
given a latency for each access model of 50 milliseconds, 2 milliseconds,
50 microseconds and 40 microseconds respectively. A strategy needs at most
1 millisecond and trades 5 million shares a month: which model, at what
monthly cost, and what does the latency requirement cost it compared with
the unconstrained choice?
\end{exercise}

\begin{exercise}[$\star\star\star$]\label{exo:m1:exchanges-brokers-venues:8}
\emph{Find the flaw.} A start-up's plan says: ``Incumbent \omterm{def:m1:exchanges-brokers-venues:exchange}{exchanges} charge
0.30 cent a share to take liquidity. We will charge 0.05 cent. Traders
minimise cost, so volume will move to us.'' Identify the quantity the plan
has left out of the trader's cost, write the trader's actual comparison, and
explain why the cheaper venue can be the more expensive one.
\end{exercise}

\section{Problem: Building an Exchange}

\begin{problem}[{Weekend problem --- the business plan of a new stock venue}]\label{pb:m1:exchanges-brokers-venues:1}
A consortium of \omterm{def:m1:exchanges-brokers-venues:broker}{brokers} plans a new venue for a stock market that trades 8
billion shares a day at an average price of \$45, over 252 days.

\medskip\noindent\textbf{Part I --- Revenue.}
\begin{enumerate}
\item The venue charges takers 0.28 cent a share and pays makers 0.25 cent.
What net fee $\nu$ does it keep per share matched?
\item Express $\nu$ in basis points of the value traded.
\item Give the annual transaction revenue at 100\% market share, and per
percentage point of share.
\item Market-data revenue would be \$60 million at 100\% share and is
proportional to share. Give total revenue per point of share.
\end{enumerate}

\medskip\noindent\textbf{Part II --- Costs and break-even.}
\begin{enumerate}[resume]
\item Technology costs \$18 million a year, regulation and surveillance \$9
million, staff \$11 million. Give the break-even share $x^\star$.
\item Express the break-even in shares a day and in dollars traded a day.
\item The consortium's members together execute 12\% of the market's volume,
and could route a third of their orders to the new venue where a match is
available. Assume matches are available for half of those. What share does
that give?
\item Is the venue profitable on members' flow alone? By how much?
\end{enumerate}

\medskip\noindent\textbf{Part III --- Buying share.}
\begin{enumerate}[resume]
\item To attract \omterm{def:m1:what-a-trading-firm-does:market-maker}{market makers} the venue pays, for its first year, a rebate
of 0.32 cent while still charging 0.28. What is $\nu$ now, and what does a
4\% share cost in transaction losses?
\item Add data revenue and fixed costs: what is the first-year loss at 4\%?
\item In year two it returns to the normal schedule and holds 6\%. Give the
profit, and the number of such years needed to repay the first-year loss.
\item An incumbent responds by cutting its own net fee, and the entrant must
match: $\nu = 0.02$ cent. Recompute $x^\star$.
\end{enumerate}

\medskip\noindent\textbf{Part IV --- Judgement.}
\begin{enumerate}[resume]
\item Why do the \omterm{def:m1:exchanges-brokers-venues:broker}{brokers} want this venue even if it only breaks even?
\item The venue considers being an ATS instead of an \omterm{def:m1:exchanges-brokers-venues:exchange}{exchange}. Name two
things it gives up and one thing it saves.
\item It considers a delay of a fraction of a millisecond on incoming orders.
Which traders does this attract, and which does it repel?
\item Explain why market share in matching has the economics of a network,
and what that implies for the number of surviving venues.
\item An \omterm{def:m1:exchanges-brokers-venues:exchange}{exchange}'s data revenue depends on its share of \emph{quotes} as
well as trades. Why might that reward posting orders that are unlikely to
execute?
\item Sponsored-access clients connect directly to the venue. What must the
venue itself verify about their sponsoring \omterm{def:m1:exchanges-brokers-venues:broker}{brokers}?
\item State the \emph{named result}: the break-even traded value per day, in
billions of dollars, under the normal fee schedule.
\item In one sentence: is this a technology business or a regulatory one?
\end{enumerate}
\end{problem}

\section{Interview questions}

\begin{interviewq}[$\star$ \iqroles{trader, developer, researcher}]\label{iq:m1:exchanges-brokers-venues:1}
Name the ways an \omterm{def:m1:exchanges-brokers-venues:exchange}{exchange} makes money. Which of them matter most to a
high-frequency trading firm's cost base?
\end{interviewq}

\begin{interviewq}[$\star$ \iqroles{developer, trader}]\label{iq:m1:exchanges-brokers-venues:2}
What is the difference between \omterm{def:m1:exchanges-brokers-venues:dma}{direct market access} and \omterm{def:m1:exchanges-brokers-venues:dma}{sponsored access}?
Who is responsible if a sponsored client's algorithm sends ten thousand
erroneous orders?
\end{interviewq}

\begin{interviewq}[$\star\star$ \iqroles{developer}]\label{iq:m1:exchanges-brokers-venues:3}
You are asked to implement pre-trade risk checks that every order must pass.
List the checks and say what constraint their latency budget puts on the
design.
\end{interviewq}

\begin{interviewq}[$\star\star$ \iqroles{researcher, trader}]\label{iq:m1:exchanges-brokers-venues:4}
Why do futures \omterm{def:m1:exchanges-brokers-venues:exchange}{exchanges} have much higher margins than stock \omterm{def:m1:exchanges-brokers-venues:exchange}{exchanges}?
\end{interviewq}

\begin{interviewq}[$\star\star$ \iqroles{trader, researcher}]\label{iq:m1:exchanges-brokers-venues:5}
A new venue offers to pay you 0.32 cent a share for posting liquidity while
the established venue pays 0.20. Your fill rate there is a fifth of what it
is on the established venue. How do you decide where to quote?
\end{interviewq}

\begin{interviewq}[$\star\star\star$ \iqroles{researcher, bank}]\label{iq:m1:exchanges-brokers-venues:6}
\omterm{def:m1:exchanges-brokers-venues:exchange}{Exchanges} sell faster data feeds and colocation to some participants and a
slower consolidated feed to everyone else. Make the strongest argument that
this is fair, then the strongest argument that it is not.
\end{interviewq}
